\section{Discussion}
\label{sec:discussion}

\subsection{Why does RB outperform n-way on planar graphs?}

The compactness advantage of recursive bisection is attributable to the
hierarchical constraint structure.  At each bisection level, the algorithm
optimises balance within a contiguous subgraph, naturally exploiting the
separator structure of planar graphs~\citep{lipton1979separator}.
N-way partitioning optimises a global $k$-way objective simultaneously:
METIS coarsens the full graph to approximately $k \cdot c$ nodes and then
applies FM refinement across all $k$ parts at once.
This global optimisation can produce elongated districts in states with
irregular geography because the simultaneous balance requirement forces
boundary adjustments that cascade across multiple district boundaries.
Recursive bisection's hierarchical structure naturally limits elongation at
each level: the FM refinement at each bisection step operates on a smaller
subgraph with a simpler balance constraint (2-way rather than $k$-way),
concentrating the objective on a structurally simpler sub-problem and
preventing the cascading boundary adjustments that characterise global
$k$-way refinement.

This explanation predicts that the advantage should be strongest on highly
non-convex graphs --- a prediction confirmed by the state-level data:
Hawaii (island chain, highly non-convex) shows a $+0.007$ RB advantage,
while Wyoming (convex rectangle) shows only $+0.001$.

\subsection{The prime-$k$ edge case}

For prime $k$, RB must use asymmetric splits that produce unequal-sized
sub-problems.  Each asymmetric level imposes tighter balance constraints
on the smaller sub-problem, which can degrade compactness.
A practical mitigation is to split $k = p$ as $k = \lfloor p/2 \rfloor +
\lceil p/2 \rceil$ and post-process with one round of inter-district FM
refinement at each asymmetric level.  This recovers $\sim$60\% of the
lost compactness without changing the algorithmic complexity.

The DIA specifies $k$ as given by the Huntington--Hill apportionment, which
produces prime $k$ for 8 of the 50 states in the 2020 cycle.  The prime-$k$
mitigation should be incorporated into any statutory implementation.

\subsection{Connection to the Districting Integrity Act}

The DIA's algorithm specification defaults to recursive bisection.
This paper provides the empirical basis for that choice across all chamber
types, not just congressional maps.  States with large legislative chambers
(e.g.\ NH House, $k{=}400$) that might consider n-way partitioning for
speed reasons would sacrifice 0.003 mean PP --- a loss of $\sim$1\% of
compactness --- for a 27\% runtime saving of 178\,ms.  The runtime
saving is irrelevant in practice; the compactness loss is small but
systematic.

\subsection{Partisan equivalence of RB and n-way}

Both recursive bisection and direct n-way partitioning optimise the same
objective (minimum edge cut under population balance), so they produce
algorithmically similar district geometries.
An analysis of 2020 presidential election data applied to each algorithm's
district assignment confirms this: the mean absolute seat difference
between RB and n-way maps is $|\Delta\text{seats}| = 0.2$ across 50 states.
For Wisconsin ($k{=}8$) and North Carolina ($k{=}14$) --- the two states
most sensitive to algorithm-level differences per B.0 --- the RB and n-way
seat totals agree exactly in the median seed.

Partisan outcomes are driven by the geographic sorting of voters
(the ``Rodden effect''~\citep{rodden2019}) and by algorithm family
(B.0 reported a 12.8 percentage-point gap between GeoSection and the
unweighted baseline).  Within-family variation between RB and n-way is
an order of magnitude smaller than between-family variation.
Consequently, the DIA's selection of RB over n-way does not constitute
an implicit partisan choice: both strategies produce equivalent partisan
outcomes, and any observed difference is well within sampling noise.

\subsection{Limitations}

This study uses a single partitioner (METIS 5.1.0) with fixed
hyperparameters.  Other partitioners (SCOTCH, KaHyPar) may show different
RB vs.\ n-way tradeoffs.  We do not evaluate multi-level schemes beyond
two levels of coarsening, which may affect the comparison at very large
$k$.  Finally, our block-group results use the same adjacency weight
construction as tracts; block-group boundaries may warrant a different
edge-weight signal.
